% ---------------------------------------------------------------------------
\section{Implementation}
\label{sec:implementation}

\sys\ is an open-source JavaScript library (MIT licence)
\todo{repository is private; make it public or deposit an archived snapshot
(e.g.\ Zenodo DOI) before submission, and cite that}. The core is a set of
pure functions over plain objects with no dependency on a map, the DOM or a
framework; it returns geometry, which a canvas renderer paints and a MapLibre
adapter positions on a map. The seven stages of \cref{sec:model} are the API:
a lens is configured by one object per stage, and every figure in this paper is
a single call with different arguments.

\paragraph{The core has no map in it.} The gallery from which
\cref{fig:gallery} is drawn runs the whole pipeline onto bare canvases with no
map library on the page. It doubles as a standing test that the core has not
acquired a rendering dependency. All raster figures in this paper are
regenerated from the library by a script in the repository
(\texttt{paper/scripts/figures.mjs}), which serves the repository, renders the
gallery and a figure page in headless Chromium, and saves each canvas.

\paragraph{Placement.} The necklace solver of \cref{sec:placement} is
independent of the rest of the library and exported on its own. It tries every
cut of the cycle and solves each exactly: by weighted pool-adjacent-violators,
or by its bounded variant when marks have feasible intervals. That makes
order-preserving placement exact, as argued there. An implementation of Speckmann and
Verbeek's algorithms is part of CartoCrow, a C++ framework for cartographic
algorithms from the Applied Geometric Algorithms group at TU
Eindhoven~\cite{cartocrow}. Its project page lists no paper for the framework
itself, so we cite the page. We found no earlier JavaScript implementation. On
2026-09-25 we searched the npm registry for \emph{necklace map},
\emph{necklace maps}, \emph{necklace}, \emph{Speckmann}, \emph{boundary
labeling} and \emph{ring map}. We also ran web and code searches for necklace
maps with JavaScript, d3, npm and GitHub, and a search of Observable. The
only npm package named \texttt{necklace} is an unrelated pipeline library.
This is a search result, not a proof of absence.

\paragraph{Fields.} A field calls the single-lens pipeline once per lattice
centre and changes nothing else. It adds two things a single lens does not
need: a lattice (hexagonal, square, triangular, or a Lloyd-relaxed centroidal
Voronoi tessellation inside a boundary, built on a Delaunay triangulation) and a
uniform grid index so that $m$ lenses over $n$ members do not cost $O(nm)$
distance tests.

Of the evaluation strategies used for toolkit research---demonstration,
usage, technical benchmarks and heuristics~\cite{ledo2018toolkits}---this paper
uses the first and third: the gallery and the coding of
\cref{sec:descriptive,sec:generative} are demonstrations, and
\cref{tab:timing} is a benchmark. A usage study with developers is not part of
this work.

\paragraph{Tests and performance.} The repository has 200 unit tests (Node's
built-in runner, no runtime dependencies), covering among other things
non-overlap, stability and overflow of the solver. \Cref{tab:timing} reports
median wall-clock times for the pure pipeline, without rendering, on the bundled
extract of 1{,}449 OpenStreetMap places.

\begin{table}[t]
  \centering
  \small
  \caption{Median time of the pure pipeline (no rendering), Node~22 on one
  core of a 2.8\,GHz Intel Xeon, 1{,}449 places; produced by
  \texttt{paper/scripts/bench.mjs}. Field sizes are the actual lattice counts.}
  \label{tab:timing}
  \begin{tabular}{@{}lr@{}}
    \toprule
    Configuration & ms \\
    \midrule
    One lens, 8 bearing sectors, necklace placement & 0.9 \\
    One lens, 24 bearing sectors & 1.2 \\
    One lens, 72 bearing sectors & 3.3 \\
    One lens, 8 sectors $\times$ 4 categories, stacked & 0.9 \\
    Field of 37 lenses, 8 sectors each & 5.2 \\
    Field of 241 lenses & 14.1 \\
    Field of 889 lenses & 17.1 \\
    \midrule
    Solver alone, 24 marks, 80\% fill & 0.4 \\
    Solver alone, 72 marks & 1.3 \\
    Solver alone, 288 marks & 17.1 \\
    \bottomrule
  \end{tabular}
\end{table}

The field of 889 lenses is only slightly slower than the field of 241 because 804 of its
lenses (and 165 of the 241) hold fewer than the minimum of three members and are
skipped after the index query, so the cost is dominated by the lenses actually
drawn (85 and 76 respectively). The solver
is $O(K^2)$ in the number of marks, which is immaterial at the bin counts a
lens uses and visible at 288.
